\section{Compute budget}
\label{sec:compute}

\subsection{Machine and measurements}

Planning is done for Tom's laptop (Apple M4 Pro: 10 performance and 4 efficiency cores, 48\,GB) until a
school compute node is confirmed. \dec{2026-10-05}

\begin{whybox}[Why budget compute first]
The advisor's instruction was that compute is the binding constraint and that architecture choices are
settled only once the budget is calculated, not guessed. Planning on the laptop is the conservative case: if
the design fits there, a school node only adds slack. Everything was measured, not assumed: a synthetic
benchmark of every configuration, then one timed run on the real pipeline to measure the overhead.
\end{whybox}

\begin{itemize}
\item \textbf{Cost of one run.} With a fixed step budget, a run costs about
  \[\text{steps}\times\text{batch}\times(6P-2\,d\,w_1)\ \text{FLOPs},\]
  where $P$ is the number of parameters, $d$ the input dimension and $w_1$ the first width (the first
  layer needs no gradient with respect to its inputs). It does not depend on how many training rows
  there are.
\item \textbf{One thread per run}: a single core reaches 30--230 GFLOP/s; for $K=3$, widths 64--32, batch
  4,096: 238 steps/s. More threads in one run give nothing.
\item \textbf{Parallel runs}: 10 runs on the performance cores give 6.9$\times$ (7.6$\times$ with the
  efficiency cores). Plan: 10 parallel processes.
\item \textbf{GPU (MPS)}: a fixed 1.6--3\,ms per step; useful only for the largest models.
\item \textbf{Real pipeline} (timed 6 October on the 2009 pilot fold, 1.13M training rows): 5.57\,ms per
  step against the benchmark's 4.19\,ms, an overhead factor of 1.33. Gate fit 29\,s (Hamilton, $K=3$,
  24 starts), base fit 0.2\,s, peak memory 3.9\,GB.
\end{itemize}

\subsection{What the design costs}

One run trains one model for one fold from one seed. The design multiplies:

{\small
\begin{tabularx}{\textwidth}{@{}X X r@{}}
\toprule
Block & Count & Runs\\
\midrule
Gate $\times$ depth grid & 6 arms (5 gates + no-regime control) $\times$ 10 seeds $\times$ 15 folds & 900 per depth\\
Primary cell, extra seeds & 20 seeds $\times$ 15 folds & 300\\
Construction arms & 5 arms $\times$ 10 seeds $\times$ 15 folds & 750\\
Base depth sweep & 10 seeds $\times$ 15 folds & 150 per depth\\
Pilot & about 50 settings $\times$ 3 folds $\times$ 3 seeds & 450\\
No-decay arm for every gate & $6\times10\times15$ & 900\\
Leave-one-episode-out & primary cell & about 30\\
\bottomrule
\end{tabularx}}

Depths 1--3: 5,580 clean runs; budgeted twice for reruns. Laptop time with 10 parallel runs:

{\small
\begin{tabular}{@{}lllrrr@{}}
\toprule
Depths & $K$ & First width & 3k steps & 10k steps & 30k steps\\
\midrule
1--3 & 3 & 64 & 7\,h & 23\,h & 71\,h\\
1--3 & 3 & 128 & 14\,h & 46\,h & 137\,h\\
1--3 & 4 & 64 & 9\,h & 30\,h & 91\,h\\
1--3 & 4 & 256 & 38\,h & 126\,h & 377\,h\\
1--5 & 4 & 256 & 58\,h & 194\,h & 581\,h\\
\bottomrule
\end{tabular}}

\dec{2026-10-05} \textbf{Envelope: depths 1--3, first width at most 128, $K$ at most 4.} Within it the whole
study, run twice, takes about 1--2 days at 10,000 steps per run and 3--6 days at 30,000, plus up to half
a day of gate fits. Compute is therefore not the binding limit inside the envelope; the statistical
limit is (Q17, Q18). Still to confirm: steps per run (from the pilot), the fit times of the jump,
Wasserstein and TVTP gates, the industry-level gate fits (Q28), and thermal slowdown over multi-day
runs.

\begin{whybox}[Why this envelope]
Inside it, the whole design (every gate, depth, seed and fold, run twice) takes days, not months, so
compute stops being the constraint and the remaining choices can be made on statistical grounds (Q17,
Q18). Outside it, the only configurations ruled out are deep (4--5 layers) and wide (256) networks with
long training, which would take about a month and which the literature gives no reason to expect to help:
Gu, Kelly \& Xiu find shallow networks suffice. Counting every arm twice allows for bugs and reruns, which
in practice double the clean count.
\end{whybox}

